% !TEX root = ../main.tex

% This chapter is the complete working manuscript for the eighth seminar
% talk. The final section contains supplementary material outside the default
% live route.

\chapter{Local derived geometry and Kuranishi charts}
\label[chapter]{chap:Local_Derived_Geometry_Kuranishi_Charts}

Talk 7 ended with the local normal-form slogan
\[
  \text{quasi-smooth}
  \quad\Longleftrightarrow\quad
  \text{locally a derived zero locus}.
\]
The right-hand side is geometric, but the word ``locally'' has not yet been
explained. Our affine description by animated $\Cinfty$-rings provides
pullbacks, Koszul presentations, and cotangent complexes. It does not visibly
display an underlying topological space, open subsets, or restriction of
functions.

The first task of this chapter is therefore to make affine derived objects
local. The basic operation is smooth localization. For a manifold $M$ and a
function $a\in\Cinfty(M)$, inverting $a$ restricts $M$ to the open set where
$a$ does not vanish. For an animated $\Cinfty$-ring $A$, the spectrum
construction organizes all such localizations into a topological space with
a sheaf of animated functions. Its guiding slogan is
\begin{quote}
  The topology of $\Spec A$ is controlled by $\pi_0A$, while its structure
  sheaf retains the full animated algebra $A$.
\end{quote}

The second task is to understand the nonuniqueness of the local zero-locus
presentations produced in this way. A triple
\[
  (U,E,s),
\]
consisting of a manifold, a vector bundle, and a section, is a presentation
of a derived zero locus. It is not the intrinsic derived object itself. We
will see this concretely by comparing
\[
  x\longmapsto x^2
  \qquad\text{and}\qquad
  (x,y)\longmapsto(x^2,y).
\]
The second presentation has one additional variable and one additional
equation, but presents the same derived object. This operation is called
stabilization.

Here $\Spec A$ is a geometric spectrum. It is unrelated to the stable
homotopy-theoretic coefficient objects mentioned in Talk 7. Stable modules
linearize an animated algebra; the geometric spectrum equips that algebra
with points, open restrictions, and a structure sheaf.

\section{From localization to derived spectra}
\label[section]{sec:Localization_To_Derived_Spectra}

\subsection{Open restriction in smooth algebra}
\label[subsection]{sec:Open_Restriction_Smooth_Algebra_Talk_Eight}

Let $M$ be a manifold, let $a\in\Cinfty(M)$, and write
\[
  D(a)=\{x\in M\mid a(x)\neq0\}.
\]
Talk 3 proved that restriction to $D(a)$ is the smooth localization at $a$.

\begin{proposition}[Smooth localization]
  \label[proposition]{prop:Smooth_Localization_Recall_Talk_Eight}
  Restriction induces an isomorphism of $\Cinfty$-rings
  \[
    \Cinfty(M)[a^{-1}]
    \iso
    \Cinfty(D(a)).
  \]
\end{proposition}

The graph
\[
  D(a)\longrightarrow M\times\R,
  \qquad
  x\longmapsto(x,a(x)^{-1})
\]
identifies $D(a)$ with the regular zero locus of
\[
  (x,y)\longmapsto ya(x)-1.
\]
This is the geometric reason that adjoining an inverse to $a$ produces
precisely the functions on $D(a)$.

It matters that this is localization in $\Cinfty$-rings. Algebraically
inverting the coordinate $x$ in the underlying commutative ring of
$\Cinfty(\R)$ produces only fractions $f(x)/x^n$. The smooth localization is
$\Cinfty(\R\setminus\{0\})$ and contains arbitrary smooth functions on the
punctured line, including functions whose growth near zero is not bounded by
any power of $x^{-1}$. The full smooth functional calculus is part of the
localization.

Every open subset of a manifold is a nonvanishing locus. If $V\subseteq M$
is open, there is a smooth function $a\colon M\to[0,\infty)$ such that
$V=D(a)$. Thus ordinary open restriction is completely visible in smooth
algebra.

\subsection{Derived localization}
\label[subsection]{sec:Derived_Localization_Talk_Eight}

Let $A$ be an animated $\Cinfty$-ring and let $a\in\pi_0A$. A morphism
$A\to B$ carries $a$ to a well-defined element of $\pi_0B$, so it makes
sense to require that image to be invertible.

\begin{definition}[Derived localization]
  \label[definition]{def:Derived_Localization_Talk_Eight}
  The \emph{localization of $A$ at $a$} is a morphism
  \[
    A\longrightarrow A[a^{-1}]
  \]
  which is initial among morphisms from $A$ carrying $a$ to an invertible
  element of $\pi_0$.
\end{definition}

This initial object is unique in the homotopical sense: its universal
property includes the homotopies and all their higher compatibilities. It is
computed by the pushout
\[
\begin{tikzcd}
  \Cinfty(\R) \rar \dar \drar[pushout]
    & A \dar \\
  \Cinfty(\R\setminus\{0\}) \rar
    & A[a^{-1}],
\end{tikzcd}
\]
where the upper morphism classifies $a$.

The decisive calculation says that this construction only restricts the
derived information already present.

\begin{theorem}[Homotopy groups of a smooth localization]
  \label[theorem]{thm:Homotopy_Groups_Smooth_Localization_Talk_Eight}
  Let $A$ be an animated $\Cinfty$-ring and let $a\in\pi_0A$. Then
  \[
    \pi_0(A[a^{-1}])
    \cong
    (\pi_0A)[a^{-1}],
  \]
  and for every $n\geq0$ the canonical morphism
  \[
    \pi_n(A)\otimes_{\pi_0A}(\pi_0A)[a^{-1}]
    \longrightarrow
    \pi_n(A[a^{-1}])
  \]
  is an isomorphism.
\end{theorem}

The theorem is Steffens, Proposition~4.1.3.13
\cite[Proposition~4.1.3.13]{Steffens_Derived_Differential_Geometry}. We quote
its proof, which uses flatness of smooth open restriction. The formula says
that localization is \emph{strong}: it creates no new higher homotopy groups.
Each existing homotopy group is simply restricted to the chosen open set.

There is also a linear version of the same assertion. Smooth localization is
formally etale, so
\[
  \mathbb L_{A[a^{-1}]/A}\simeq0.
\]
Transitivity therefore gives
\begin{equation}
  \label{eq:Cotangent_Complex_Localizes_Talk_Eight}
  \mathbb L_{A[a^{-1}]}
  \simeq
  \mathbb L_A\otimes_AA[a^{-1}].
\end{equation}
This is the bridge back to Talk 7: the stable linearization of an open
restriction is the restriction of the stable linearization.

If $A$ is homotopically finitely presented, then so is $A[a^{-1}]$. This
follows either from the displayed pushout or from Steffens,
Remark~4.1.3.14
\cite[Remark~4.1.3.14]{Steffens_Derived_Differential_Geometry}.

\subsection{Points, basic opens, and animated functions}
\label[subsection]{sec:Derived_Spectrum_Slogan_Talk_Eight}

We now organize the localizations of $A$ into a geometric object. Assume for
the moment that $A$ is homotopically finitely presented. Its set of real
points is
\[
  |\Spec A|
  =
  \Hom_{\Cinfty\text{-rings}}(\pi_0A,\R).
\]
For $a\in\pi_0A$, define the basic open
\[
  D(a)
  =
  \{x\colon\pi_0A\to\R\mid x(a)\neq0\}.
\]
These opens generate a topology. On $D(a)$, place the full animated
localization $A[a^{-1}]$ rather than only its discrete truncation.

\begin{construction}[Derived spectrum]
  \label[construction]{con:Derived_Spectrum_Talk_Eight}
  The \emph{derived spectrum} of $A$ is the locally animated
  $\Cinfty$-ringed space
  \[
    \Spec A
    =
    \bigl(|\Spec A|,\mathcal O_{\Spec A}\bigr),
  \]
  whose structure sheaf is obtained from the assignments
  \[
    \mathcal O_{\Spec A}(D(a))
    \simeq
    A[a^{-1}].
  \]
\end{construction}

For the finite-presentation objects used in the live argument, these
assignments already have the expected locality properties. The ordinary
structured-space construction and its extension to animated rings are
discussed in \Cref{sec:Supplement_Real_Spectrum_Talk_Eight}.

Truncation has an especially simple description:
\[
  t_0(\Spec A)
  =
  \bigl(|\Spec A|,\pi_0\mathcal O_{\Spec A}\bigr)
  \cong
  \Spec(\pi_0A).
\]
The topological space does not change. Truncation only discards the higher
homotopy sheaves of functions.

A one-point example makes the distinction visible. Let
\[
  K(0)=\bigl(\R\otimes\Lambda(e),\partial e=0\bigr)
\]
be the animated function algebra of the derived self-intersection
$*\times_{\R}*$. Then
\[
  \pi_0K(0)=\R,
  \qquad
  \pi_1K(0)=\R.
\]
The following three objects have the same one-point underlying space:
\[
\begin{array}{@{}lll@{}}
  \text{object} & \pi_0\text{ of the stalk} & \pi_1\text{ of the stalk}\\
  \hline
  \text{ordinary point} & \R & 0\\
  \text{dual-number point} & \R[\varepsilon]/(\varepsilon^2) & 0\\
  \text{derived self-intersection} & \R & \R.
\end{array}
\]
Topology distinguishes none of them. The ordinary and higher parts of the
structure sheaf distinguish all three.

\subsection{The affine comparison theorem}
\label[subsection]{sec:Affine_Comparison_Talk_Eight}

There is always a natural morphism
\[
  A\longrightarrow\Gamma(\Spec A,\mathcal O_{\Spec A}).
\]
For arbitrary animated $\Cinfty$-rings it need not be an equivalence. In the
finite-presentation range of derived manifolds, it is.

\begin{definition}[Affine derived $\Cinfty$-scheme]
  \label[definition]{def:Affine_Derived_Cinfty_Scheme_Talk_Eight}
  An \emph{affine derived $\Cinfty$-scheme of finite presentation} is a
  locally animated $\Cinfty$-ringed space equivalent to $\Spec A$ for a
  homotopically finitely presented animated $\Cinfty$-ring $A$.
\end{definition}

\begin{theorem}[Affine comparison in finite presentation]
  \label[theorem]{thm:Affine_Comparison_Finite_Presentation_Talk_Eight}
  Spectrum and global sections induce an equivalence of $\infty$-categories
  \[
    \operatorname{Alg}_{\Cinfty}(\An)_{\mathrm{fp}}\catop
    \simeq
    \mathrm{d}\Cinfty\mathrm{Aff}_{\mathrm{fp}}.
  \]
  In particular, for every homotopically finitely presented animated
  $\Cinfty$-ring $A$, the unit induces an isomorphism
  \[
    A
    \iso
    \Gamma(\Spec A,\mathcal O_{\Spec A}).
  \]
\end{theorem}

The proof passes through the class of fair animated $\Cinfty$-rings. Steffens
first proves spectrum and global sections to be inverse on fair objects and
then proves that almost finitely presented animated $\Cinfty$-rings are fair.
We use only the resulting finite-presentation statement. The relevant results
are Theorem~4.1.3.22, Proposition~4.1.3.33, and Corollary~4.1.3.34
\cite[Theorem~4.1.3.22, Proposition~4.1.3.33, and
Corollary~4.1.3.34]{Steffens_Derived_Differential_Geometry}.

The theorem does not introduce a second affine theory. It equips the affine
objects of Talks 4 to 7 with an explicit topology, open restrictions, and a
structure sheaf. Combining it with the algebraic presentation from Talk 4
gives
\[
  \mathrm{DMfd}
  \simeq
  \mathrm{d}\Cinfty\mathrm{Aff}_{\mathrm{fp}}.
\]

\section{Local derived smooth geometry}
\label[section]{sec:Local_Derived_Smooth_Geometry_Talk_Eight}

The affine comparison theorem gives a geometric meaning to a single animated
algebra. We now allow these affine objects to be restricted and glued.

\subsection{Open subobjects and locally affine schemes}
\label[subsection]{sec:Open_Subobjects_Locally_Affine_Talk_Eight}

Let $(X,\mathcal O_X)$ be a locally animated $\Cinfty$-ringed space and let
$V\subseteq|X|$ be open. Restriction gives another such space
\[
  (V,\mathcal O_X|_V)
\]
and a morphism to $X$. We call this an \emph{open subobject}.

On an affine derived spectrum, principal open restriction is precisely the
geometric form of localization:
\begin{equation}
  \label{eq:Principal_Open_Derived_Spectrum}
  \Spec A[a^{-1}]
  \simeq
  \bigl(D(a),\mathcal O_{\Spec A}|_{D(a)}\bigr).
\end{equation}
The underlying topological spaces agree by construction. On every smaller
principal open, the structure sheaves agree by the universal properties of
the corresponding iterated localizations.

\begin{definition}[Derived $\Cinfty$-scheme]
  \label[definition]{def:Locally_Affine_Derived_Cinfty_Scheme_Talk_Eight}
  A \emph{derived $\Cinfty$-scheme locally of finite presentation} is a
  locally animated $\Cinfty$-ringed space which admits an open cover by
  affine derived $\Cinfty$-schemes of finite presentation.
\end{definition}

This is a genuine enlargement of the affine category. The universal
$\infty$-category $\mathrm{DMfd}$ used so far is equivalent to the affine
finite-presentation category. A locally affine derived $\Cinfty$-scheme may
instead require several affine opens.

The sheaf condition is what permits these affine restrictions to be glued.
For the live route, it is enough to remember its familiar meaning: animated
functions defined locally and agreeing coherently on all intersections arise
from one function on the union. The Cech formulation and the distinction
between sheaf descent here and moduli descent in Talk 9 are recorded in the
supplement.

\subsection{Truncation and quasi-smoothness are local}
\label[subsection]{sec:Truncation_Quasi_Smoothness_Local_Talk_Eight}

The ordinary truncation of a derived $\Cinfty$-scheme is
\[
  t_0X
  =
  \bigl(|X|,\pi_0\mathcal O_X\bigr).
\]
On an affine open it restricts to
\[
  t_0(\Spec A)=\Spec(\pi_0A).
\]
Again the topological space is unchanged. Only the animated structure sheaf
has been truncated.

The cotangent complex restricts to opens by
\Cref{eq:Cotangent_Complex_Localizes_Talk_Eight}. This lets us globalize the
finiteness condition from Talk 7.

\begin{definition}[Quasi-smooth derived $\Cinfty$-scheme]
  \label[definition]{def:Quasi_Smooth_Derived_Cinfty_Scheme_Talk_Eight}
  A derived $\Cinfty$-scheme locally of finite presentation is
  \emph{quasi-smooth} if, on an affine open cover $\Spec A_i$, every
  cotangent complex $\mathbb L_{A_i}$ is perfect of Tor-amplitude contained
  in $[-1,0]$.
\end{definition}

\begin{proposition}[Locality of quasi-smoothness]
  \label[proposition]{prop:Locality_Quasi_Smoothness_Talk_Eight}
  The condition in
  \Cref{def:Quasi_Smooth_Derived_Cinfty_Scheme_Talk_Eight} is independent of
  the chosen affine open cover and may be checked after any affine open
  refinement.
\end{proposition}

\begin{proof}
  \Cref{eq:Cotangent_Complex_Localizes_Talk_Eight} shows that
  perfection and the stated Tor-amplitude persist under restriction.
  Conversely, perfection and Tor-amplitude are local properties of stable
  modules on an affine spectrum. If they hold on an open cover, the local
  duals and amplitude bounds agree on overlaps and give the corresponding
  global properties.
\end{proof}

We can now restate the normal-form theorem from Talk 7 without leaving the
word ``locally'' undefined.

\begin{theorem}[Local zero-locus form]
  \label[theorem]{thm:Local_Zero_Locus_Form_Talk_Eight}
  Let $X$ be a derived $\Cinfty$-scheme locally of finite presentation. Then
  $X$ is quasi-smooth if and only if every point of $t_0X$ has an open
  neighborhood equivalent to the derived zero locus of a section of a
  finite-rank vector bundle over a manifold.
\end{theorem}

One direction was calculated in Talk 7: the cotangent complex of a
finite-rank zero locus is perfect in degrees $-1$ and $0$. The converse is a
derived inverse-function and normal-form theorem. We quote it from Steffens,
Proposition~5.1.1.11 and Corollary~5.1.1.13
\cite[Proposition~5.1.1.11 and
Corollary~5.1.1.13]{Steffens_Derived_Differential_Geometry}.

The converse is a substantial input. The tangent-complex calculation alone
does not produce local coordinates or equations. It only verifies the
amplitude of an object already presented as a zero locus.

\subsection{One global zero locus and its local equations}
\label[subsection]{sec:Global_Zero_Locus_Local_Charts_Talk_Eight}

Let $q\colon E\to M$ be a finite-rank vector bundle and let
$s\colon M\to E$ be a section. The global derived zero locus
\[
  Z=Z^{\mathrm{der}}(s)
\]
was constructed before any trivialization of $E$ was chosen:
\[
\begin{tikzcd}
  Z \rar \dar \drar[pullback]
    & M \dar["s"] \\
  M \rar["0"'] & E.
\end{tikzcd}
\]

Choose a bundle-trivializing open cover $M=\bigcup_iU_i$ and identify
\[
  E|_{U_i}\cong U_i\times\R^r.
\]
In this frame the section is a smooth map
\[
  s_i=(s_{i,1},\ldots,s_{i,r})
  \colon
  U_i\longrightarrow\R^r.
\]
Let $Z_i\subseteq Z$ be the inverse image of $U_i$ along the projection
$Z\to M$.

\begin{proposition}[Local Koszul presentations]
  \label[proposition]{prop:Local_Koszul_Presentations_Talk_Eight}
  There are natural equivalences
  \[
    Z_i
    \simeq
    \Spec K(s_i),
  \]
  where
  \[
    K(s_i)
    =
    \bigl(
      \Cinfty(U_i)\otimes\Lambda(e_1,\ldots,e_r),
      \partial e_j=s_{i,j}
    \bigr).
  \]
\end{proposition}

\begin{proof}
  Open restriction commutes with the pullback defining $Z$, so
  \[
    Z_i
    \simeq
    U_i\times_{E|_{U_i}}U_i.
  \]
  After the chosen trivialization, this is the derived zero locus of
  $s_i\colon U_i\to\R^r$. The Koszul calculation from Talk 5 identifies its
  animated function algebra with $K(s_i)$.
\end{proof}

On an overlap $U_i\cap U_j$, the two local formulas are related by the frame
change of $E$. More intrinsically, both formulas present the restriction of
the same global pullback $Z$. Their comparison is therefore not additional
structure imposed on $Z$. It is induced by restriction from the one object
which was present before either frame was chosen.

The tangent complexes behave in the same way. Globally along the classical
zero set,
\[
  \mathbb T_Z
  \simeq
  [TM|_{Z(s)}\xrightarrow{Ds}E|_{Z(s)}]
\]
in cohomological degrees $0$ and $1$. In a frame on $U_i$, this restricts to
\[
  [TU_i|_{Z(s_i)}\xrightarrow{Ds_i}U_i\times\R^r].
\]
The change of frame acts simultaneously on the obstruction bundle and on the
derivative. The local complexes are presentations of one global complex.

This example joins the two halves of the chapter. The spectrum construction
explains how one global derived zero locus acquires local equations. We now
ask how much of one such system of equations belongs to the object itself.

\section{Kuranishi charts as presentations}
\label[section]{sec:Kuranishi_Charts_As_Presentations_Talk_Eight}

\subsection{Presentation, object, and chart}
\label[subsection]{sec:Presentation_Object_Chart_Talk_Eight}

We work throughout without isotropy groups, boundaries, or corners. This is
the finite-dimensional local situation needed for the later elliptic
applications.

\begin{definition}[Kuranishi presentation]
  \label[definition]{def:Kuranishi_Presentation_Talk_Eight}
  A \emph{Kuranishi presentation without isotropy} is a triple
  \[
    K=(U,E,s),
  \]
  where $U$ is a manifold, $E\to U$ is a finite-rank vector bundle, and
  $s\colon U\to E$ is a section. It presents the derived zero locus
  \[
    \mathbf Z^{\mathrm{der}}(K)
    :=
    Z^{\mathrm{der}}(s).
  \]
  Its ordinary zero locus is
  \[
    Z(K)=\{x\in U\mid s(x)=0_x\}.
  \]
\end{definition}

When a presentation is used as a chart for a fixed derived
$\Cinfty$-scheme $X$, the data also include an equivalence
\[
  X|_V
  \simeq
  \mathbf Z^{\mathrm{der}}(K)
\]
for some open $V\subseteq|X|$. On ordinary truncations this gives a
homeomorphism from $V$ to the zero set $Z(K)$, often called the footprint
map. In informal usage the triple itself is frequently called a Kuranishi
chart. The distinction between presentation and chart relative to $X$ is
useful whenever choices are being compared.

There are therefore three levels:

\begin{enumerate}
  \item The \emph{presentation} $(U,E,s)$ consists of auxiliary
    finite-dimensional data.
  \item The \emph{presented object} is the intrinsic derived zero locus
    $\mathbf Z^{\mathrm{der}}(K)$.
  \item A \emph{comparison} is a morphism of presentations which induces a
    morphism between the presented derived objects.
\end{enumerate}

At $x\in Z(K)$, Talk 7 gives
\begin{equation}
  \label{eq:Tangent_Complex_Kuranishi_Presentation_Talk_Eight}
  \mathbb T_x\mathbf Z^{\mathrm{der}}(K)
  \simeq
  [T_xU\xrightarrow{D_xs}E_x]
\end{equation}
in cohomological degrees $0$ and $1$. Thus
\[
  \operatorname{Def}_x(K)=\ker(D_xs),
  \qquad
  \operatorname{Obs}_x(K)=\operatorname{coker}(D_xs),
\]
and
\[
  \operatorname{vdim}K
  =
  \dim U-\operatorname{rank}E.
\]

The triple contains more information than the object it presents. It records
one ambient manifold, one obstruction bundle, and one system of equations.
The remainder of the chapter exhibits operations which change these choices.

\subsection{The tangent complex is not a complete invariant}
\label[subsection]{sec:Tangent_Complex_Not_Complete_Talk_Eight}

Before comparing presentations, we must rule out a tempting shortcut.
Consider
\[
  f(x)=x^2
  \qquad\text{and}\qquad
  g(x)=x^3.
\]
Both derivatives vanish at the origin, so their tangent complexes agree:
\[
  \mathbb T_0Z^{\mathrm{der}}(f)
  \simeq
  [\R\xrightarrow{0}\R]
  \simeq
  \mathbb T_0Z^{\mathrm{der}}(g).
\]
Their ordinary truncations are nevertheless different. Talk 5 gives
\[
  \pi_0\mathcal O\bigl(Z^{\mathrm{der}}(f)\bigr)
  \cong
  \frac{\Cinfty(\R)}{(x^2)},
  \qquad
  \pi_0\mathcal O\bigl(Z^{\mathrm{der}}(g)\bigr)
  \cong
  \frac{\Cinfty(\R)}{(x^3)}.
\]
Repeated use of Hadamard's lemma identifies the corresponding local rings at
the origin with
\[
  \R[\varepsilon]/(\varepsilon^2)
  \qquad\text{and}\qquad
  \R[\varepsilon]/(\varepsilon^3).
\]
They are not isomorphic.

\begin{warning}[Linear data do not classify the nonlinear object]
  \label[warning]{warn:Tangent_Complex_Not_Classifier_Talk_Eight}
  Isomorphic deformation and obstruction spaces, or even isomorphic abstract
  tangent complexes, do not imply that two derived zero loci are equivalent.
  A comparison of nonlinear presentations requires a morphism between them.
\end{warning}

This is analogous to the ordinary inverse-function theorem. The derivative
of a specified smooth map may prove that the map is locally invertible. An
accidental isomorphism between two tangent spaces does not produce a map of
manifolds. Later we will quote a derived version of this principle for a
specified morphism of Kuranishi presentations.

\subsection{Shrinking and strict coordinate changes}
\label[subsection]{sec:Elementary_Changes_Kuranishi_Talk_Eight}

Two changes of presentation are immediate.

First, let $W\subseteq U$ be open. Restriction gives
\[
  K|_W=(W,E|_W,s|_W).
\]
Its derived zero locus is the corresponding open restriction:
\begin{equation}
  \label{eq:Kuranishi_Shrinking_Open_Restriction}
  \mathbf Z^{\mathrm{der}}(K|_W)
  \simeq
  \mathbf Z^{\mathrm{der}}(K)|_{W\cap Z(s)}.
\end{equation}
If $W$ contains all of $Z(s)$, the ambient presentation has been shrunk
without changing the complete presented object. If $W$ contains only part of
$Z(s)$, the result presents a proper open subobject. The word ``shrinking''
is used for both operations, so the footprint should always be stated.

Second, suppose that $f_b\colon U\to V$ is a diffeomorphism and
$f_v\colon E\to F$ is a vector-bundle isomorphism over $f_b$ satisfying
\[
  f_v\circ s=t\circ f_b.
\]
Then $(f_b,f_v)$ is a strict coordinate and frame change. It carries the
pullback square defining $Z^{\mathrm{der}}(s)$ isomorphically to the pullback
square defining $Z^{\mathrm{der}}(t)$.

These operations change the visible equations only mildly. Stabilization is
the first operation which changes both the dimension of the ambient manifold
and the rank of the obstruction bundle while preserving the derived object.

\section{Stabilization and invariance}
\label[section]{sec:Stabilization_And_Invariance_Talk_Eight}

\subsection{Adding a variable and a transverse equation}
\label[subsection]{sec:Definition_Stabilization_Talk_Eight}

Let $K=(U,E,s)$ be a Kuranishi presentation and let $p\colon G\to U$ be a
finite-rank vector bundle. On its total space, the pullback bundle $p^*G$ has
a tautological section
\[
  \tau_G\colon\operatorname{Tot}(G)\longrightarrow p^*G,
  \qquad
  \tau_G(g)=g\in G_{p(g)}.
\]

\begin{definition}[Stabilization]
  \label[definition]{def:Stabilization_Kuranishi_Talk_Eight}
  The \emph{stabilization of $K$ by $G$} is the Kuranishi presentation
  \[
    K^G
    =
    \left(
      \operatorname{Tot}(G),
      p^*E\oplus p^*G,
      (p^*s,\tau_G)
    \right).
  \]
\end{definition}

The new equation $\tau_G=0$ forces the new ambient variable to lie on the
zero section $U\subseteq\operatorname{Tot}(G)$. Since this equation is
transverse, it should contribute no residual derived structure. The original
equation $s=0$ then remains.

Projection $p\colon\operatorname{Tot}(G)\to U$ together with projection
onto $p^*E$ defines a morphism of presentations
\begin{equation}
  \label{eq:Stabilization_Morphism_Talk_Eight}
  \sigma_G\colon K^G\longrightarrow K.
\end{equation}
Indeed, projection carries $(p^*s,\tau_G)$ to $s\circ p$.

\subsection{The stabilized square equation}
\label[subsection]{sec:Stabilized_Square_Equation_Talk_Eight}

Begin with
\[
  K
  =
  (\R,\R\times\R,x\mapsto x^2).
\]
Stabilization by the trivial line gives
\[
  K^{\R}
  =
  (\R^2,\R^2,(x,y)\mapsto(x^2,y)).
\]
We compare the two presentations at three levels.

At the level of ordinary zero loci, both equations have only the origin as a
solution. Projection $(x,y)\mapsto x$ identifies the two zero sets.

At the level of stable linearization, the original tangent complex is
\[
  \mathbb T_0\mathbf Z^{\mathrm{der}}(K)
  \simeq
  [\R\xrightarrow{0}\R].
\]
The stabilized tangent complex is
\[
  \mathbb T_0\mathbf Z^{\mathrm{der}}(K^{\R})
  \simeq
  \left[
    \R^2
    \xrightarrow{
      \left(
        \begin{smallmatrix}
          0&0\\
          0&1
        \end{smallmatrix}
      \right)}
    \R^2
  \right].
\]
It decomposes as
\begin{equation}
  \label{eq:Square_Stabilization_Tangent_Decomposition}
  \left[
    \R^2
    \xrightarrow{
      \left(
        \begin{smallmatrix}
          0&0\\
          0&1
        \end{smallmatrix}
      \right)}
    \R^2
  \right]
  \cong
  [\R\xrightarrow{0}\R]
  \oplus
  [\R\xrightarrow{1}\R].
\end{equation}
The second summand is zero in the stable module category. Projection induces
a quasi-isomorphism from the stabilized complex to the original complex.

The complete nonlinear calculation takes place on animated function
algebras. The two Koszul presentations are
\[
  \mathcal K(K)
  =
  \bigl(
    \Cinfty(\R)\otimes\Lambda(e),
    \partial e=x^2
  \bigr)
\]
and
\[
  \mathcal K(K^{\R})
  =
  \bigl(
    \Cinfty(\R^2)\otimes\Lambda(e,f),
    \partial e=x^2,
    \partial f=y
  \bigr).
\]
The pair $(y,f)$ is the Koszul resolution of the transverse equation $y=0$.
Eliminating it restricts smooth functions to the $x$-axis and induces an
equivalence
\begin{equation}
  \label{eq:Square_Stabilization_Koszul_Equivalence}
  \mathcal K(K^{\R})
  \simeq
  \mathcal K(K).
\end{equation}

We have checked stabilization at three levels:

\begin{enumerate}
  \item The ordinary zero loci are identified.
  \item The tangent complexes differ by an acyclic summand.
  \item The full animated function algebras are equivalent.
\end{enumerate}

The third statement is the intrinsic conclusion. The first two are its
ordinary and infinitesimal shadows. This is different from the comparison of
$x^2$ and $x^3$: there we had an accidental agreement of tangent complexes
but no equivalence of the nonlinear algebras.

\subsection{Why stabilization preserves the derived object}
\label[subsection]{sec:General_Stabilization_Proof_Talk_Eight}

The preceding calculation globalizes without choosing a frame for $G$.

\begin{theorem}[Invariance under stabilization]
  \label[theorem]{thm:Stabilization_Preserves_Derived_Object_Talk_Eight}
  For every Kuranishi presentation $K=(U,E,s)$ and every finite-rank vector
  bundle $G\to U$, projection induces an equivalence
  \[
    \mathbf Z^{\mathrm{der}}(K^G)
    \simeq
    \mathbf Z^{\mathrm{der}}(K).
  \]
\end{theorem}

\begin{proof}
  The tautological section $\tau_G$ is transverse to the zero section of
  $p^*G$. Its ordinary zero locus is the zero section
  \[
    U\longrightarrow\operatorname{Tot}(G).
  \]
  By transverse exactness, its derived zero locus is this ordinary manifold
  $U$, with no additional derived structure.

  The derived zero locus of the direct-sum section
  $(p^*s,\tau_G)$ can be formed in two stages. First impose
  $\tau_G=0$, and then impose the restriction of $p^*s=0$. Finite pullbacks
  may be pasted in either order, so this iterated zero locus is equivalent to
  the zero locus of the pair. After the first stage the ambient manifold is
  the zero section $U$, the bundle $p^*E$ restricts to $E$, and the section
  $p^*s$ restricts to $s$. Hence the second stage is precisely
  $Z^{\mathrm{der}}(s)$.
\end{proof}

The tangent calculation is compatible with this equivalence. Along the zero
section there is a canonical splitting
\[
  T_{0_x}\operatorname{Tot}(G)
  \cong
  T_xU\oplus G_x.
\]
The derivative of the stabilized section is block diagonal:
\[
  \left[
    T_xU\oplus G_x
    \xrightarrow{
      \left(
        \begin{smallmatrix}
          D_xs&0\\
          0&1
        \end{smallmatrix}
      \right)}
    E_x\oplus G_x
  \right].
\]
Consequently,
\begin{equation}
  \label{eq:General_Stabilization_Tangent_Decomposition}
  \mathbb T_x\mathbf Z^{\mathrm{der}}(K^G)
  \simeq
  \mathbb T_x\mathbf Z^{\mathrm{der}}(K)
  \oplus
  [G_x\xrightarrow{1}G_x].
\end{equation}
The virtual dimension is unchanged:
\[
  \dim\operatorname{Tot}(G)
  -\operatorname{rank}(p^*E\oplus p^*G)
  =
  \dim U-\operatorname{rank}E.
\]

\subsection{Recognizing safe changes of presentation}
\label[subsection]{sec:Recognizing_Safe_Changes_Talk_Eight}

Stabilization was proved directly. There is also a general practical
criterion for a specified comparison map.

\begin{definition}[Morphisms of Kuranishi presentations]
  \label[definition]{def:Morphism_Kuranishi_Presentation_Talk_Eight}
  Let $K=(U,E,s)$ and $L=(V,F,t)$. A \emph{morphism of presentations}
  \[
    f=(f_b,f_v)\colon K\longrightarrow L
  \]
  consists of a smooth map $f_b\colon U\to V$ and a fiberwise-linear map
  $f_v\colon E\to F$ over $f_b$ such that
  \[
    f_v\circ s=t\circ f_b.
  \]
\end{definition}

The compatibility carries $Z(s)$ into $Z(t)$ and induces a morphism
\[
  \mathbf Z^{\mathrm{der}}(f)
  \colon
  \mathbf Z^{\mathrm{der}}(K)
  \longrightarrow
  \mathbf Z^{\mathrm{der}}(L).
\]
At a zero $x$, differentiation gives a map of tangent complexes
\begin{equation}
  \label{eq:Tangent_Map_Kuranishi_Morphism_Talk_Eight}
  \begin{tikzcd}
    T_xU \rar["D_xs"] \dar["D_xf_b"']
      & E_x \dar["(f_v)_x"] \\
    T_{f_b(x)}V \rar["D_{f_b(x)}t"']
      & F_{f_b(x)}.
  \end{tikzcd}
\end{equation}

\begin{definition}[Weak equivalence of presentations]
  \label[definition]{def:Weak_Equivalence_Kuranishi_Talk_Eight}
  A morphism $f\colon K\to L$ is a \emph{weak equivalence} if it induces a
  bijection $Z(K)\to Z(L)$ and the map in
  \Cref{eq:Tangent_Map_Kuranishi_Morphism_Talk_Eight} is a
  quasi-isomorphism at every zero.
\end{definition}

\begin{theorem}[Presentation-invariance criterion]
  \label[theorem]{thm:Weak_Equivalence_Presents_Equivalence_Talk_Eight}
  A morphism of finite-dimensional Kuranishi presentations is a weak
  equivalence if and only if its induced morphism of derived zero loci is an
  equivalence.
\end{theorem}

The direction from derived equivalence to weak equivalence follows by taking
ordinary truncations and tangent complexes. The converse is a derived
inverse-function theorem. Pointwise quasi-isomorphism makes the relative
cotangent complex vanish at every zero. Finite presentation and derived
Nakayama extend this vanishing to open neighborhoods, where the morphism is
etale. The bijection of zero loci then globalizes these local equivalences.
A fuller source-based proof is recorded in
\Cref{sec:Supplement_Weak_Equivalence_Kuranishi}.

The theorem must not be used without its input map. The comparison of $x^2$
and $x^3$ supplies isomorphic abstract tangent complexes but no morphism whose
linearization is being tested. The theorem recognizes equivalences; it does
not construct nonlinear comparisons from linear data.

\section{Summary and supplementary materials}
\label[section]{sec:Summary_And_Supplementary_Materials_Talk_Eight}

\subsection{Takeaways and the coherence problem}
\label[subsection]{sec:Takeaways_Coherence_Problem_Talk_Eight}

Let us collect the route of the chapter.
\begin{enumerate}
  \item Smooth localization is open restriction written on functions.
  \item The topology of $\Spec A$ is controlled by $\pi_0A$, while its
    structure sheaf retains the full animated $\Cinfty$-ring $A$.
  \item In finite presentation, affine derived spectra and animated
    $\Cinfty$-rings contain the same information.
  \item Quasi-smoothness is local because cotangent complexes restrict under
    smooth localization.
  \item Quasi-smooth derived $\Cinfty$-schemes are locally derived zero loci.
  \item A Kuranishi presentation is auxiliary data presenting such a local
    derived object.
  \item Stabilization changes the presentation by an acyclic variable-equation
    pair and preserves the presented derived object.
  \item Tangent complexes test a specified comparison map, but do not classify
    nonlinear derived objects by themselves.
\end{enumerate}

There is one important problem which this chapter has deliberately not
solved. Suppose that a moduli problem is covered by local Kuranishi
presentations $K_i$. On a double overlap one wants a comparison between
$K_i$ and $K_j$. On a triple overlap, the two composites of comparisons need
not agree strictly, so one needs a homotopy between them. On a quadruple
overlap, those homotopies must themselves be compatible. The pattern
continues.

Pairwise weak equivalences therefore do not constitute a global Kuranishi
object. The next chapter replaces a list of chosen charts by an intrinsic
moduli functor and packages all overlap data in an $\infty$-categorical
descent condition. Its main representability input will be the following
result.

\begin{theorem}[Open local representability, preview]
  \label[theorem]{thm:Open_Local_Representability_Preview}
  Let $\mathcal X$ be a derived $\Cinfty$-stack with an effective open atlas
  by affine derived $\Cinfty$-schemes belonging to a class stable under open
  restriction. Then $\mathcal X$ is represented by a derived
  $\Cinfty$-scheme locally belonging to that class.
\end{theorem}

The definitions and proof belong to Talk~9. For now, the point is only the
shape of the answer: local presentations glue when their comparison data are
the restrictions of one object satisfying descent. This is why the local
geometry developed here is an input to, rather than a substitute for, the
coherence theorem.

\subsection{Supplement: the real spectrum and its structure sheaf}
\label[subsection]{sec:Supplement_Real_Spectrum_Talk_Eight}

The live talk used the derived spectrum at slogan level. We now record the
ordinary construction more carefully. It explains where the topology and
local rings come from.

\begin{definition}[Real spectrum]
  \label[definition]{def:Real_Spectrum_Talk_Eight}
  Let $A$ be a $\Cinfty$-ring. Its \emph{real spectrum} is the set
  \[
    \Spec_{\R}A
    =\Hom_{\Cinfty\text{-}\mathrm{Ring}}(A,\R).
  \]
  For $a\in A$, set
  \[
    D(a)=\{x\in\Spec_{\R}A\mid x(a)\neq0\}.
  \]
  The subsets $D(a)$ form a basis for a topology on $\Spec_{\R}A$.
\end{definition}

Indeed,
\[
  D(1)=\Spec_{\R}A,
  \qquad
  D(a)\cap D(b)=D(ab).
\]
Every $a\in A$ defines an evaluation function
\[
  a^*\colon\Spec_{\R}A\longrightarrow\R,
  \qquad
  a^*(x)=x(a).
\]
The spectrum topology is the weakest topology for which all evaluation
functions are continuous. To see the nontrivial direction, choose, for an
open $V\subseteq\R$, a smooth function $h$ with
$V=\{t\mid h(t)\neq0\}$. Then
\[
  (a^*)^{-1}(V)=D(\Phi_h(a)).
\]
It is also Hausdorff: two distinct homomorphisms differ on some $a\in A$, and
disjoint intervals around the two values pull back to disjoint
neighborhoods.

If $A=\Cinfty(\R^n)/I$, evaluation identifies the real spectrum with
\[
  Z(I)=\{p\in\R^n\mid f(p)=0\text{ for all }f\in I\}
\]
with its ordinary subspace topology. Thus the real spectrum recovers the
expected point set, but it cannot see nilpotent or higher derived structure.

For $x\in\Spec_{\R}A$, define the local $\Cinfty$-ring
\[
  A_x=A[(A\setminus x^{-1}(0))^{-1}].
\]
Its elements are germs of elements of $A$ near $x$. The structure sheaf may
be described without making a choice of basis presentation.

\begin{construction}[Structure sheaf]
  \label[construction]{con:Structure_Sheaf_Talk_Eight}
  For $U\subseteq\Spec_{\R}A$ open, let $\mathcal O_A(U)$ be the set of
  families
  \[
    \sigma(x)\in A_x,
    \qquad x\in U,
  \]
  which are locally represented by elements of $A$. Pointwise smooth
  operations and restriction make $\mathcal O_A$ a sheaf of
  $\Cinfty$-rings. Define
  \[
    \Spec A=(\Spec_{\R}A,\mathcal O_A).
  \]
\end{construction}

By construction, $\mathcal O_{A,x}\cong A_x$. Moreover, localization induces
an isomorphism of local $\Cinfty$-ringed spaces
\begin{equation}
  \label{eq:Principal_Open_Ordinary_Spectrum_Talk_Eight}
  \Spec A[a^{-1}]
  \iso
  (D(a),\mathcal O_A|_{D(a)}).
\end{equation}
For a finitely presented $\Cinfty$-ring, the basis assignment
$D(a)\mapsto A[a^{-1}]$ already gives the same sheaf after sheafification.
For arbitrary rings, the local-germ construction avoids subtleties caused by
different localization maps presenting the same real open subset.

\begin{theorem}[Reconstruction of manifolds]
  \label[theorem]{thm:Reconstruction_Manifolds_Talk_Eight}
  For a manifold $M$, evaluation induces an isomorphism of local
  $\Cinfty$-ringed spaces
  \[
    \Spec\Cinfty(M)\iso(M,\Cinfty_M).
  \]
\end{theorem}

\begin{proof}
  Every homomorphism $\Cinfty(M)\to\R$ is evaluation at a unique point of
  $M$. The basic open $D(a)$ is the nonvanishing locus of the smooth function
  $a$. Conversely, if $x\in U\subseteq M$ with $U$ open, a bump function
  supported in $U$ and nonzero at $x$ produces a basic neighborhood
  $x\in D(a)\subseteq U$. Hence the spectrum topology is the manifold
  topology.

  Smooth localization identifies the stalk at $x$ with the ring
  $\Cinfty_{M,x}$ of smooth germs. Finally, a section of $\mathcal O_A$ is
  locally represented by global smooth functions and is therefore smooth.
  Conversely, cutoff functions show that every smooth function on an open
  subset is locally the restriction of a global smooth function.
\end{proof}

For an animated $\Cinfty$-ring, the same construction is applied to
$\pi_0A$ to obtain the point set and topology, while the localizations
$A[a^{-1}]$ provide the animated structure sheaf. The homotopy groups of
that sheaf are the sheafifications of
\[
  D(a)\longmapsto
  \pi_n(A)\otimes_{\pi_0A}(\pi_0A)[a^{-1}].
\]
This makes precise the slogan that the topology sees $\pi_0A$, whereas the
structure sheaf sees all of $A$.

\subsection{Supplement: the Möbius self-intersection}
\label[subsection]{sec:Supplement_Mobius_Self_Intersection_Talk_Eight}

The zero section of a nontrivial bundle shows why local equations need not
admit a global choice of generators. Let $L\to S^1$ be the Möbius line bundle
and form the derived self-intersection of its zero section:
\[
  X=S^1\times_LS^1.
\]
Its classical truncation is $S^1$.

Choose an open cover $S^1=U_0\cup U_1$ on which $L$ is trivial. On $U_i$,
the zero section is represented by the zero function, and hence
\[
  \mathcal O_X|_{U_i}
  \simeq
  \bigl(\Cinfty_{U_i}\otimes\Lambda(e_i),\partial e_i=0\bigr).
\]
Thus $\pi_0\mathcal O_X|_{U_i}\cong\Cinfty_{U_i}$ and
$\pi_1\mathcal O_X|_{U_i}\cong\Cinfty_{U_i}e_i$.

Let $g_{01}\colon U_0\cap U_1\to\R^\times$ be the transition function of
$L$. The degree-one Koszul generator transforms as a local frame of
$L^\vee$:
\[
  e_1=g_{01}^{-1}e_0.
\]
Consequently,
\[
  \pi_1\mathcal O_X
  \cong
  \bigl(U\longmapsto\Gamma(U,L^\vee)\bigr).
\]
There is no global nowhere-vanishing degree-one generator. Dually, the
tangent complex is
\[
  \mathbb T_X\simeq[TS^1\xrightarrow{0}L],
\]
so its obstruction sheaf is the Möbius bundle itself. The derived object is
global, but every elementary equation displaying it is local and depends on
a frame.

\subsection{Supplement: weak equivalences of presentations}
\label[subsection]{sec:Supplement_Weak_Equivalence_Kuranishi}

We explain the source of
\Cref{thm:Weak_Equivalence_Presents_Equivalence_Talk_Eight}. Let
$f\colon K\to L$ be a morphism of the finite-dimensional presentations
considered in this chapter, and write
\[
  F=\mathbf Z^{\mathrm{der}}(f)
  \colon X\longrightarrow Y
\]
for the induced morphism of finite-presentation affine derived
$\Cinfty$-schemes.

If $F$ is an equivalence, its truncation is an isomorphism and its tangent
map is a quasi-isomorphism at every point. Thus $f$ is a weak equivalence.
For the converse, suppose that $f$ induces a bijection on zero loci and a
quasi-isomorphism on tangent complexes. The fiber of the relative cotangent
complex $\mathbb L_{X/Y}$ at a point $x$ is dual to the cofiber of the map of
tangent complexes at $x$. It therefore vanishes. Since $\mathbb L_{X/Y}$ is
perfect in the present finite-presentation setting, derived Nakayama implies
that it vanishes on a neighborhood of every point. Hence $F$ is etale near
every point.

The derived inverse-function theorem then identifies a neighborhood of $x$
with a neighborhood of $F(x)$. Because the map on underlying zero loci is a
bijection, these local equivalences cover both source and target and assemble
to an equivalence $X\simeq Y$. The finite-presentation and geometricity
hypotheses are essential here. The statement is not a principle saying that
arbitrary derived objects with isomorphic tangent complexes are equivalent.

There is also an algebraic way to recognize what is happening. After
trivializing the bundles, the presentations produce Koszul algebras. A
morphism of presentations induces a morphism between these algebras. Weak
equivalence says that the induced map has the same classical points and is a
quasi-isomorphism on the linearized complexes at those points. The argument
above promotes that pointwise linear information, through the relative
cotangent complex and finite presentation, to a local equivalence of the
full nonlinear Koszul algebras.

\subsection{Supplement: terminology and truncations}
\label[subsection]{sec:Supplement_Terminology_Truncations_Talk_Eight}

The terminology surrounding Kuranishi geometry is not uniform. In this
chapter a \emph{Kuranishi presentation} means only a finite-dimensional
triple $(U,E,s)$ without isotropy, boundary, or corners. A \emph{Kuranishi
chart for $X$} additionally includes an identification of its derived zero
locus with an open subobject of $X$. This usage isolates exactly the local
calculation needed later.

Other theories package different amounts of information. Joyce's
d-manifolds and d-orbifolds are 2-categorical geometric structures, while
Kuranishi spaces are built from charts and coordinate changes with their own
specified coherence. Fully derived manifolds retain higher mapping animae
which cannot in general be recovered from a 2-categorical truncation. These
frameworks are closely related, but they should not be identified without a
comparison theorem and its hypotheses. In particular, the local
stabilization calculation in this chapter does not by itself compare their
global notions of atlas.

\subsection{Related literature}
\label[subsection]{sec:Related_Literature_Talk_Eight}

The derived localization theorem, affine comparison, local zero-locus normal
form, and weak-equivalence criterion used in the chapter come from Steffens
\cite[Proposition~4.1.3.13, Corollary~4.1.3.34, Corollary~5.1.1.13, and
Remarks~1.0.0.4 and~1.0.0.6]{Steffens_Derived_Differential_Geometry}.
Joyce gives detailed accounts of ordinary $\Cinfty$-spectra, affine
$\Cinfty$-schemes, and the reconstruction of manifolds
\cite[Sections~2.2 and~3.1]{Joyce_Introduction_Cinfty_Schemes}
\cite[Sections~4.4 and~4.5]{Joyce_Algebraic_Geometry_Cinfty_Rings}.
The universal derived-manifold viewpoint and its relation to local
zero-locus constructions are developed by Carchedi and Steffens
\cite{Carchedi_Steffens_Universal_Derived_Manifolds}. The representability
theorem previewed at the end of the live route is
\cite[Proposition~2.1.48]{Steffens_Representability_Elliptic_Moduli}.
